\section{Answer to Reviewer rM5r}


\subsection{Answer to W1}

\textit{The LKH integration result is not very convincing. LKH-3's native candidates use alpha-nearness — itself a 1-tree/HK construction — so C2TSP produces candidates of the type LKH already expects, and its comparatively better results are natural.}

\textbf{\underline{RE:}}
\textbf{Are the better LKH integration results natural?}
No. The structural alignment alone is not sufficient. The advantage comes from the quality of the learned signal, which captures solution structure within the connected rooted 1-tree family. 
% and the gain comes from the content of the learned signal. 
Two observations show that the advantage does not follow automatically from this alignment.
 
First, vanilla LKH-3 (H0) already uses $\alpha$-nearness candidates computed by the classical Held–Karp subgradient procedure. If producing ``candidates of the type LKH expects'' were sufficient, C2TSP integration could at best match H0. Instead, replacing LKH's internally computed $\alpha$-nearness with C2TSP's learned candidates and duals ($H2\rightarrow H4$) improves on vanilla LKH-3 for $n\leq 500$ at comparable wall time, and stays within $\sim 1\text{\textpertenthousand}$ of it at every size. The experiment therefore tests signal quality: C2TSP recovers 1-tree structure accurate enough to outperform the solver's own hand-tuned construction of the same object.
 
% Second, NeuroLKH is a direct counterexample to the ``natural advantage'' reading. It is trained specifically to produce the inputs LKH expects (supervised Held–Karp node penalties and candidate sets), yet it degrades by two orders of magnitude out of distribution. Producing the expected type of signal is thus clearly not sufficient for good integration, what matters is whether the learned signal is correct, structural, and size-consistent.

Second, and most directly, the additional baselines in our response to Q1 hold the $H2$ budget fixed: non-learned $\alpha$-nearness gives gaps (\text{\textpertenthousand}) of 5.93-6.23, versus 0.11-0.20 for C2TSP. The learned edge ranking therefore adds information beyond structural compatibility.

\subsection{Answer to W2}

\textit{All experiments are on uniformly sampled 2D Euclidean TSP, with training at TSP100 and zero-shot testing up to TSP2000. There is no TSPLIB evaluation, no non-uniform or structured distributions. This limits the generality of the representation claims.}

\textbf{\underline{RE:}}
\textbf{Test on TSPLIB, structured distribution.}
We thank the reviewer for this point and have added a TSPLIB benchmark$^{[1]}$ covering both structured node distributions and non-Euclidean metrics. The results show that the performance advantage extends beyond uniformly sampled Euclidean instances. Full per-instance results will be included in the appendix of the revision.

\textbf{Setup.} All methods use the same TSP-100 trained checkpoints as in the main text, applied zero-shot with no fine-tuning. Aligned with Table~1 of the paper, we evaluate on every TSPLIB instance with fewer than 1000 nodes: 50 structured instances and 27 non-Euclidean instances.

\begin{table}[h]
\centering
\caption{\textbf{Zero-shot transfer to TSPLIB.} Cells report the mean optimality gap (\%) $\pm$ standard deviation across the instances, and the per-instance wall time in seconds (in parentheses).}
\label{tab:tsplib-core}
\small
\setlength{\tabcolsep}{4pt}
\begin{tabular}{lrrrr}
\toprule
 & \multicolumn{2}{c}{Structured (50 inst.)} & \multicolumn{2}{c}{Non-Euclidean (27 inst.)} \\
\cmidrule(lr){2-3}\cmidrule(lr){4-5}
Method & L1 $\downarrow$ & L2 $\times 100$ $\downarrow$ & L1 $\downarrow$ & L2 $\times 100$ $\downarrow$ \\
\midrule
\textbf{C2TSP} & \textbf{10.61}$\pm$7.27 (3.55) & \textbf{2.46}$\pm$1.93 (3.55) & \textbf{7.36}$\pm$3.96 (2.90) & \textbf{1.81}$\pm$1.94 (2.90) \\
DIFUSCO & 47.50$\pm$29.86 (2.35) & 11.72$\pm$15.34 (2.35) & 31.98$\pm$26.96 (1.84) & 7.22$\pm$11.42 (1.84) \\
Fast-T2T & 27.72$\pm$17.16 (0.06) & 7.06$\pm$8.48 (0.07) & 25.67$\pm$17.39 (0.05) & 5.13$\pm$6.97 (0.06) \\
DIMES & 16.74$\pm$5.80 (0.02) & 3.42$\pm$1.76 (0.03) & 19.13$\pm$10.29 (0.02) & 3.95$\pm$3.66 (0.02) \\
UTSP & 36.06$\pm$15.00 (0.02) & 9.18$\pm$6.18 (0.02) & 29.09$\pm$21.96 (0.01) & 7.26$\pm$7.77 (0.01) \\
\bottomrule
\end{tabular}
\end{table}

\textbf{Results.}
The ranking from the main text transfers to TSPLIB: C2TSP attains the best mean gap in all cells.

\textbf{Explanation.}
% We attribute this pattern to the division of labor in learning architecture. Heatmap-based baselines ask the network to produce the solution structure directly. In C2TSP, learning is restricted to a residual edge field $\Gamma_\theta$ on top of the true costs $D$, while the global structure is supplied by fixed algorithmic layers, including exact rooted 1-tree inference, HK equilibration, and certificate-guided sharpening. The guarantees of these layers are proved for arbitrary node distributions and distance metrics. 
Heatmap-based baselines learn edge scores directly from the training distribution. In C2TSP, the network only learns a residual edge field, while connectivity and expected degree balance are supplied by the rooted 1-tree family and HK equilibration. Because these guarantees are independent of the node distribution and distance metric, they reduce the amount of global structure the network must learn and improve zero-shot transfer.

\subsection{Answer to W3}

\textit{Training wall-clock cost is not reported. Each forward pass involves iterative sharpening pass and HK pass, so training cost is non-obvious and should be stated for a fair comparison with baselines.}

\textbf{\underline{RE:}}
\textbf{Report the training wall time.}
We thank the reviewer for this suggestion, and we will add a full training-cost accounting to the appendix of the revision. All models were trained for 100 epochs on the same machine on the same TSP-100 dataset:

\begin{table}[h]
\centering
\small
\caption{Model size and training time comparison.}
\label{tab:model_size}
\begin{tabular}{lrrr}
\toprule
Method & Params & Train inst. & Wall time \\
\midrule
C2TSP    & 123,750   & 2,048 & 1 h 16 m 40 s \\
DIFUSCO  & 5,333,762 & 2,048 & 52 m 05 s \\
Fast-T2T & 5,333,762 & 2,048 & 1 h 16 m 12 s \\
DIMES    & 64,321    & 2,048 & 15 m 25 s \\
UTSP     & 36,072    & 2,048 & 1 m 51 s \\
\bottomrule
\end{tabular}
\end{table}

 % | Method   | Params    | Train inst.       | Wall time     |
 % |----------|-----------|-------------------|---------------|
 % | C2TSP    | 123,750   | 2048              | 1 h 16 m 40 s |
 % | DIFUSCO  | 5,333,762 | 2048              | 52 m 05 s     |
 % | Fast-T2T | 5,333,762 | 2048              | 1 h 16 m 12 s |
 % | DIMES    | 64,321    | 2048              | 15 m 25 s     |
 % | UTSP     | 36,072    | 2048              | 1 m 51 s      |

We summarize here three major findings. (1) The reviewer's observation is correct, and the table quantifies it: each C2TSP forward pass runs one HK equilibration plus $K$ sharpening re-equilibrations , and this dominates its training cost. (2) The absolute cost nonetheless stays within the same range as the baselines: essentially identical to Fast-T2T and $\sim 1.5\times$ DIFUSCO on identical data. (3) The comparison slightly favors DIFUSCO and Fast-T2T: DIFUSCO and Fast-T2T are supervised and additionally require pre-generation of optimal tours as labels for training, and this cost is not included in the wall time. It is small at $n=100$ but grows with instance size (exponentially), whereas C2TSP trains entirely label-free.

\subsection{Answer to W4}

\textit{The root $r$ is fixed but TSP is node-symmetric, and sensitivity to root choice is never discussed.}

\textbf{\underline{RE:}}
\textbf{Is performance sensitive to the choice of the fixed root?} We performed additional experiments and verified that the performance is not affected by the fixed root node. Specifically, we evaluated the trained model on the TSP-50 testset 50 times, once with each node selected as the root: the gaps range from 1.528\% to 1.717\% with a standard deviation of 0.046\%, so performance is stable across all root choices.

% Two design choices already limit root sensitivity. First, all nodes in each training graph are i.i.d. samples, so their indices have no inherent meaning. Second, we randomly relabel node indices during training, so different original nodes occupy the root position. The evaluation below tests sensitivity at inference, keeping the same trained model and TSP-50 test instances as in the paper.
% The original root-0 run reproduced the reported result (1.550\%) in Table 3 of the paper.
All numbers in the table below are optimality gaps in $\%$.

\begin{table}[htbp]
    \centering
    % \caption{Summary statistics of the evaluated roots.}
    \label{tab:root_gap_statistics}
    \begin{tabular}{lrrrrr}
        \hline
        Roots evaluated & Mean & Std. & Min  & Max  & Median \\
        \hline
        Root 0--49 & 1.626 & 0.046 & 1.528 & 1.717 & 1.634 \\
        \hline
    \end{tabular}
\end{table}

% Across all root choices, the range is 0.189\%, demonstrating the stability and robustness of our model.

%  | Roots evaluated | Mean gap (%) | Sample std. (%) | Min gap (%) | Max gap (%) | Median gap (%) |
% | ---: | ---: | ---: | ---: | ---: | ---: |
% | root 0-49 | 1.625826 | 0.045795 | 1.528071 | 1.716668 | 1.634371 |


\subsection{Answer to Q1}

\textit{1. Under the same LKH candidate settings, have the authors tested a no-learned-signal control — distance-only, random $k$-NN, or non-learned 1-tree candidates? This would tell us whether the LKH robustness comes from learning or from producing candidates that are structurally compatible with LKH's alpha-nearness machinery.}

\textbf{\underline{RE:}}
\textbf{Does the LKH robustness come from learning or from structural compatibility?} From learning. We ran the requested controls. Structurally compatible but non-learned
candidates are far worse than C2TSP.

% \textbf{Setup.}
% Every control row uses a parameter file identical to C2TSP's $H2$ setting and draws its five candidates from the same 20-NN pool. Only the ranking score differs. 
% Table R3 shows the gaps (\text{\textpertenthousand}) against Concorde optimal value, 100 instances per size:

\begin{table}[h]
\centering
\caption{\textbf{Cells report gaps (\text{\textpertenthousand}) against Concorde optimal value, 100 instances per size.}}
\label{tab:lkh-control}
\small
\setlength{\tabcolsep}{6pt}
\renewcommand{\arraystretch}{1.05}
\begin{tabular}{lrrr}
\toprule
Candidate source & TSP100 & TSP200 & TSP500 \\
\midrule
Random 5 of 20-NN                            & 666.44 & 568.66 & 345.84 \\
Distance 5-NN                                &   7.30 &   7.93 &   6.86 \\
Non-learned $1$-tree ($\alpha$-nearness) & 6.04 & 6.23 & 5.93 \\
LKH-3 vanilla ($\alpha$-nearness + HK ascent) & 0.38 &   0.22 &   0.36 \\
Distance 5-NN + C2TSP initial tour           &   2.49 &   2.67 &   4.35 \\
\textbf{C2TSP (H2)}                          & \textbf{0.11} & \textbf{0.13} & \textbf{0.20} \\
\bottomrule
\end{tabular}
\end{table}

We summarize here three major findings. (1) Structural compatibility is not the explanation. Random 5-of-20-NN candidates are equally compatible with LKH's machinery, yet yield gaps of 346–666. Relative to the distance-only 5-NN ranking under identical $H2$ settings, replacing the ranking score by C2TSP's learned marginal is worth 34--66$\times$. (2) Non-learned 1-tree candidates yield gaps of 5.9–6.2. Adding Held–Karp subgradient ascent to it recovers vanilla LKH-3, so the classical ascent is worth 16--28$\times$. (3) C2TSP's initial tour on top of distance-only candidates reaches only 2.49–4.35: the learned information is primarily carried by the candidate ranking. The full table along with the discussion will be added to the appendix.

\subsection{Answer to Q2}

\textit{2. The root $r$ is fixed throughout, but the family $\mathcal{U}_r$, the HK dual, and the marginals all depend on $r$. Is performance stable across random root choices? Does averaging over multiple roots help?}

\textbf{\underline{RE:}}
\textbf{Is performance stable across roots, and does averaging over roots help?} Performance is stable across roots, and averaging does not help.
The 50-root evaluation in our response to W4 gives a gap range of 1.528\%–1.717\% with a standard deviation of 0.046\% on TSP-50. Averaging model outputs over 10 roots leaves the L1–L3 decoders essentially unchanged and clearly hurts L4 decoder.

\textbf{Setup.} We conducted an evaluation using one trained TSP50 model on 1,000 test instances. The single-root baseline uses root 0. For the 10-root ensemble, we average the model outputs before decoding, while retaining root 0 as the decoder root. The table reports mean optimality gaps in \%.

\begin{center}
\begin{tabular}{lrrrrrr}
\hline
Root setting & L1 & L2$\times$1 & L2$\times$10 & L2$\times$100 & L3 & L4 \\
\hline
Single root 0
& 6.100 & 4.229 & 1.328 & 1.316 & 2.966 & 1.920 \\
Average of roots 0--9
& 6.176 & 4.274 & 1.324 & 1.302 & 2.988 & \textbf{3.973} \\
\hline
\end{tabular}
\end{center}

\textbf{Results.} Root averaging has little effect on L1--L3 under the same decoding budget. Only L4 deteriorates from 1.920\% to 3.973\%. Since L4 selects a root-0 rooted 1-tree using perturbed edge costs, averaging costs across roots disrupts the root-dependent score structure and makes the 1-tree selection less coherent. We therefore retain single-root inference, and will add this result to the appendix.

\subsection{Answer to Q3}

\textit{3. Are forward-map constraints better than loss penalties in terms of final performance? Could the authors ablate: (a) forward HK only; (b) penalty only; (c) both — reporting tour quality and degree-balance statistics? This would clarify whether forward equilibration offers a measurable advantage over standard penalization.}

\textbf{\underline{RE:}}
\textbf{Are forward-map constraints better than loss penalties in terms of final performance?} Yes. In the requested ablation, penalty-only training gives a 20.81\% gap, forward HK alone gives 2.21\%, and adding the penalty on top of forward HK brings no further gain. We will also include this additional results to the appendix of the camera-ready version of our paper. 

\textbf{Setup.}
We disable stage~2 (edge sharpening) to isolate the effect of forward equilibration. For penalty-based variants, the training loss is
\[
L=\mathrm{cost}+5\lVert r\rVert_2^2,
\]
where \(r\) is the vector of degree-2 residuals for each graph. The table below reports \(\lVert r\rVert_2\) averaged across test graphs. We compare (a) forward HK only, (b) penalty only, and (c) forward HK with the penalty. Results use the TSP50 test set, and all gaps are reported in \%.

\begin{table}[htbp]
    \centering
    \label{tab:hk_penalty_ablation}
    \begin{tabular}{lrr}
        \hline
        Method & Gap (\%) & Mean residual \(\lVert r\rVert_2\) \\
        \hline
        (a) Forward HK only & 2.206 & 0.0341 \\
        (b) Penalty only & 20.811 & 0.6959 \\
        (c) Forward HK + penalty & 2.212 & 0.0337 \\
        \hline
    \end{tabular}
\end{table}

\textbf{Result.} Forward HK provides a clear measurable advantage over penalty-only training: removing the forward equilibration increases the gap from approximately \(2.21\%\) to \(20.81\%\) and substantially worsens the residual. In contrast, (a) and (c) achieve nearly identical tour quality and residuals, indicating that the additional penalty provides no measurable benefit once HK equilibration is imposed in the forward map.


\subsection{Answer to Q4}

\textit{4. Please add multi-seed results (mean and standard deviation) for the core tables.}

\textbf{\underline{RE:}}
\textbf{Add multi-seed results for the core tables.} We appreciate this suggestion and we repeated the full evaluation of Tables 1 and 2 of the paper with $10$ random seeds on the same test instances. No ranking in Table 1 or Table 2 flips under any seed, so all conclusions are unchanged. The full multi-seed tables will be updated in the revision.

There are two additional remarks worth further elaboration. (1) Many cells are deterministic by design: C2TSP, DIMES, and UTSP inference is deterministic, and the $L1/L2$ decoders contain no randomness. (2) Where variance exists, it is negligible relative to the reported effects: the largest seed std in the entire table is $\pm 0.28\%$ (UTSP L3, TSP500), and every C2TSP cell has a std of at most $\pm 0.07\%$.
% , as shown below:
% \begin{table}[h]
% \centering
% \small
% \caption{Multi-seed C2TSP results (Table 1, partial).}
% \begin{tabular}{lccccc}
% \toprule
% & TSP50 & TSP100 & TSP200 & TSP500 & TSP1000 \\
% \midrule
% L3 & $3.12 \pm 0.05$ & $5.09 \pm 0.02$ & $7.93 \pm 0.05$ & $13.45 \pm 0.07$ & $20.42 \pm 0.06$ \\
% L4 & $2.30 \pm 0.03$ & $4.61 \pm 0.04$ & $8.75 \pm 0.05$ & $15.51 \pm 0.02$ & $21.66 \pm 0.03$ \\
% \bottomrule
% \end{tabular}
% \end{table}


\textbf{Reference}

${[1]}$ Reinelt, Gerhard. ``TSPLIB—A traveling salesman problem library.'' ORSA journal on computing 3.4 (1991): 376-384.